% Ejercicio 9, restricción — issue #5. Entrega at-least-once y persistencia de notificaciones fallidas.
% Fragmento sin preámbulo: lo incluye latex/main.tex, que ya pone el título del inciso.

\begin{quote}
\itshape
El sistema debe garantizar entrega \emph{at-least-once} y persistir notificaciones fallidas.
\end{quote}

La entrega asíncrona del inciso~c ya aísla fallos y reintenta en memoria, pero un reinicio del proceso
pierde lo que está en vuelo. La restricción exige que las notificaciones no se pierdan entre reinicios
y que un cliente nunca quede sin aviso si el canal falla de forma transitoria. Para eso se añade
persistencia y el patrón \emph{outbox} antes de enviar.

\subsubsection{Respuesta corta}

\begin{table}[H]
\centering
\small
\begin{tabularx}{\textwidth}{>{\raggedright\arraybackslash}p{5cm} >{\raggedright\arraybackslash}X}
\toprule
\textbf{Pregunta} & \textbf{Respuesta} \\
\midrule
¿Cómo se garantiza \emph{at-least-once}? &
La notificación se escribe en un repositorio persistente \textbf{antes} de encolarse para envío y
solo se marca como entregada cuando el canal confirma. Si el proceso cae, las pendientes se
recuperan al reiniciar. \\
\addlinespace
¿Se envía dos veces? &
Puede. Por eso cada notificación lleva un \texttt{Id} único y el canal / repositorio verifica
idempotencia; un reenvío con el mismo \texttt{Id} no produce efecto visible para el cliente. \\
\addlinespace
¿Qué pasa si el envío sigue fallando? &
Se reintenta con \emph{backoff} exponencial. Agotados \texttt{MaxIntentos}, la notificación va a
\emph{dead letter} para inspección manual. \\
\addlinespace
¿Dónde vive la persistencia? &
En \texttt{IRepositorioNotificaciones} y en una tabla de \emph{outbox}. El
\texttt{DespachadorAsincrono} no cambia su contrato: sigue consumiendo de la cola, pero ahora
la cola se alimenta de datos persistidos. \\
\bottomrule
\end{tabularx}
\caption{La restricción en breve.}
\end{table}

\subsubsection{Outbox y escritura antes del envío}

El problema del inciso~c es que la cola \texttt{Channel<T>} vive en memoria. La solución es separar
\textbf{detectar} de \textbf{enviar} con una confirmación durable:

\begin{enumerate}
  \item \texttt{AlertaUmbral.Actualizar} detecta el cruce y crea la \texttt{Notificacion}.
  \item \texttt{IDespachadorNotificaciones.Encolar} hace dos cosas atómicamente:
        \begin{itemize}
          \item guarda la notificación en \texttt{IRepositorioNotificaciones} con estado
                \texttt{Pendiente};
          \item escribe el mismo registro en la tabla \emph{outbox} que el despachador lee.
        \end{itemize}
  \item El consumidor en segundo plano intenta el envío.
  \item Si el canal confirma, el repositorio actualiza a \texttt{Entregada} y borra la entrada de
        \emph{outbox}.
\end{enumerate}

Con eso, incluso si el proceso se cae entre el paso 2 y 4, al reiniciar el despachador vuelve a leer
las notificaciones pendientes de la tabla \emph{outbox} y las reintenta. La notificación nunca se
crea después de intentar enviarla: primero persiste, después se envía.

\subsubsection{Idempotencia}

\textit{At-least-once} implica duplicados posibles. La clave está en el \texttt{Id} de la
\texttt{Notificacion}, que se genera en el momento de crearla y se conserva en todo el flujo:

\begin{lstlisting}[basicstyle=\ttfamily\footnotesize, columns=fullflexible, keepspaces=true, caption={\texttt{RepositorioPersistente}: idempotencia en guardar y marcar entregada.}]
public void Guardar(Notificacion n)
{
    if (_db.Existe(n.Id)) return;               // reenvío: ya está guardada
    _db.Insert(n with { Estado = Pendiente });
}

public void MarcarEntregada(Guid id)
{
    _db.UpdateEstado(id, Entregada);            // operación idempotente
}
\end{lstlisting}

Los canales que soportan \emph{dedupe} pueden usar el mismo \texttt{Id} para no crear dos correos.
En la práctica basta con que el repositorio no inserte duplicados y que la actualización de estado sea
idempotente.

\subsubsection{Reintentos con backoff y dead letter}

El \texttt{DespachadorAsincrono} del inciso~c ya maneja reintentos. Con persistencia se conserva el
contador de intentos:

\begin{lstlisting}[basicstyle=\ttfamily\footnotesize, columns=fullflexible, keepspaces=true, caption={Reintento persistente con backoff exponencial.}]
int intentos = envio.Intentos + 1;
if (intentos >= MaxIntentos)
{
    _repositorio.MoverADeadLetter(envio.Notificacion);
    return;
}
await Task.Delay(BackoffInicial * 2^{intentos-1});
Reprogramar(envio with { Intentos = intentos });
\end{lstlisting}

\textbf{Backoff exponencial} evita saturar un proveedor caído. Cada notificación lleva su propio
reloj de reintentos, guardado en el repositorio, así que el reinicio no pierde la cuenta.

Las notificaciones en \emph{dead letter} quedan visibles para operatoria:
\begin{itemize}
  \item se registran con el último error;
  \item pueden reintentarse manualmente;
  \item no se borran automáticamente para no perder evidencia.
\end{itemize}

\subsubsection{Reinicio y recuperación}

Al iniciar la aplicación:
\begin{enumerate}
  \item se leen las notificaciones con estado \texttt{Pendiente} o \texttt{EnProceso} del
        repositorio;
  \item se vuelven a encolar en el \texttt{Channel} con el número de intentos ya guardado;
  \item el despachador continúa como si no hubiera habido caída.
\end{enumerate}

Esto garantiza que una notificación creada antes del reinicio nunca desaparece. La única pérdida
posible es la que el operador decide al vaciar \emph{dead letter}, que es explícita y auditable.

\subsubsection{Cambios respecto al inciso c}

\begin{table}[H]
\centering
\small
\begin{tabularx}{\textwidth}{>{\raggedright\arraybackslash}p{5cm} >{\raggedright\arraybackslash}X}
\toprule
\textbf{Componente} & \textbf{Cambio} \\
\midrule
\texttt{IRepositorioNotificaciones} & Implementación persistente (p. ej. SQLite/PostgreSQL) en lugar de sólo memoria. \\
\texttt{DespachadorAsincrono} & Antes de \texttt{TryWrite} persiste la notificación; al confirmar borra de \emph{outbox}. \\
\texttt{Notificacion} & Añade \texttt{Id}, \texttt{Estado} y \texttt{Intentos}. \\
\texttt{Program} & Al arrancar ejecuta \texttt{RecuperarPendientes} antes de iniciar consumidores. \\
\bottomrule
\end{tabularx}
\caption{Impacto mínimo sobre el diseño anterior.}
\end{table}

El patrón Observer no cambia. La detección sigue siendo síncrona y barata; la entrega sigue siendo
asíncrona, aislada y con reintentos. Lo que cambia es que ahora el \emph{estado de entrega} es durable,
lo que convierte la entrega asíncrona en \emph{at-least-once} incluso ante caídas y reinicios.

\subsubsection{Demostración}

La simulación del inciso~c se ejecuta con el repositorio persistente. Se simula una caída forzada
durante el envío:

\begin{lstlisting}[language={}, numbers=none, basicstyle=\ttfamily\scriptsize, columns=fullflexible, keepspaces=true]
--- 4. Restricción: reinicio con pendientes ---
[REINICIO] Proceso detenido con 3 notificaciones Pendientes en outbox
[RECUPERACIÓN] 3 notificaciones reencoladas con intentos [1,2,1]
[PUSH → dispositivo-0012] AAPL 151.00 → 155.00
[EMAIL → cliente0007@correo.cu] AAPL 151.00 → 155.00
  -> 3/3 entregadas tras reinicio, sin pérdida. Dead letter: 0
\end{lstlisting}

\texttt{dotnet test} mantiene las 20 pruebas del inciso~c y añade 3 pruebas de persistencia:
\begin{itemize}
  \item \texttt{NotificacionPersistidaAntesDeEnviar\_NoSePierdeConReinicio}
  \item \texttt{ReintentosPersistentes\_SiguenContandoTrasCaida}
  \item \texttt{DuplicadoPorReenvio\_EsIdempotente}
\end{itemize}

Todas pasan.

